\documentclass[10pt,a4paper]{amsart}
\usepackage[margin=19mm]{geometry}
\usepackage{amsmath,amssymb,mathtools}
\usepackage[scr=rsfs,cal=euler]{mathalfa}
\usepackage{xfrac}
\usepackage[unicode,hidelinks,bookmarks=false]{hyperref}
\newcommand{\ie}{i.e.\ }
\newcommand{\cf}{cf.\ }
\newcommand{\resp}{respectively\ }
\newcommand{\Subsection}{\subsection}
\providecommand{\nobreakdash}{\nobreak\hbox{-}}
\setlength{\emergencystretch}{2em}
\multlinegap0pt
\usepackage{bm}
\usepackage{bbm}
\usepackage{dsfont}
\usepackage{xspace}
\usepackage{cancel}
\usepackage{mathtools}
\usepackage{amsthm}
\makeatletter
\@ifundefined{theorem}{\newtheorem{theorem}{Theorem}}{}
\makeatother
\title[Wave-Particle Decomposition for Kinetic Equations I: Theory ]{Wave-Particle Decomposition for Kinetic Equations I: Theory and Numerics}
\author{Chang Liu, Kun Xu}
\date{Source: arXiv:2606.26915v1 (CC BY 4.0)}
\begin{document}
\maketitle

\noindent\emph{Source and licence.} Wave-Particle Decomposition for Kinetic Equations I: Theory and Numerics. Version arXiv:2606.26915v1 (CC BY 4.0), distributed under the Creative Commons Attribution 4.0 International licence, as recorded in the arXiv licence metadata for this exact version and shown on the abstract page https://arxiv.org/abs/2606.26915v1. Full rights notice: licenses/arxiv-2606.26915-CC-BY-4.0.txt.\par\smallskip

\noindent\textit{Editorial scope.} Counted unit: the Theorem whose source label is \texttt{thm:ap} in the article (source0.tex lines 918--939, source label \texttt{thm:ap}), delivered together with the article's own complete original proof of it (source0.tex lines 941--967, one \texttt{proof} environment of 27 lines). No mathematics is written, repaired or re-derived: statement and proof are the article's own text line for line. Required context retained uncounted: eq:fv\_update\_total (lines 578--590); eq:wave\_flux\_algorithm (lines 652--662); eq:algorithm\_A01 (lines 683--687). Every definition, display and label of those lines is retained unchanged. Omitted from this package and not redistributed: 1--577, 591--651, 663--682, 688--917, 940--940, 968--2206. Nothing inside a retained excerpt is omitted or altered: every retained body is delivered exactly as the article's lines are written, so no omission receipt of any kind accompanies this package. Citations: the retained excerpts carry no citation command at all, so no bibliography entry of the article is transcribed and no reference is redistributed. Reference adaptation: none was needed and none was made; every reference inside the delivered unit resolves inside it, so no unresolved reference or citation is produced. Printed slips kept literal: no printed word, symbol, hypothesis or number of the counted statement or of its proof was altered. Layout and class: the article's own class is \texttt{elsarticle} with packages including amsmath, amssymb, amsthm, geometry, setspace, booktabs, graphicx, subcaption, placeins, tikz, hyperref; the wrapper is the lane's pinned \texttt{amsart} editorial wrapper at 10pt on A4, which restates the theorem-like environments the retained text needs and adds the article's own macro definitions that the retained text needs (none), with extra wrapper packages (bm, bbm, dsfont, xspace, cancel, mathtools). The delivered numbering is the unit's own. Assets: no retained span contains a figure environment, a graphics-inclusion command, a PDF-page inclusion, a table or a TikZ picture; no image file is copied into this package, the package declares no asset dependency and no image file is redistributed with it. Licence: version arXiv:2606.26915v1 of this article is distributed under the Creative Commons Attribution 4.0 International licence, as recorded in the arXiv licence metadata for this exact version and shown on the abstract page https://arxiv.org/abs/2606.26915v1. A full rights notice is delivered at licenses/arxiv-2606.26915-CC-BY-4.0.txt. Characters: every retained excerpt is pure ASCII, so the delivered engine, XeLaTeX with its default Latin Modern text font, reports no missing character.
\begin{equation}\label{eq:fv_update_total}
    \mathbf U_i^{n+1}
    =
    \mathbf U_i^n
    -
    \frac{\Delta t}{|\Omega_i|}
    \sum_{j\in N(i)}
    \left(
    \mathbf F_{\mathcal{W},ij}
    +
    \mathbf F_{\mathcal{P},ij}
    \right)|\Gamma_{ij}|,
\end{equation}

\begin{equation}\label{eq:wave_flux_algorithm}
    \mathbf F_{\mathcal{W},ij}
    =
    \frac{1}{\Delta t}\int_0^{\Delta t}
    \left\langle (\boldsymbol{\xi}\cdot\boldsymbol{n}_{ij})\,\boldsymbol{\psi}\,f_{\mathcal{W},ij}(t,\boldsymbol{\xi}) \right\rangle\,\mathrm{d} t
    =
    \left\langle (\boldsymbol{\xi}\cdot\boldsymbol{n}_{ij})\,\boldsymbol{\psi}\,
    \mathcal{G}_{ij}^{0}\left(
    b_{0,ij}+b_{a,ij}\,a_{ij}+b_{t,ij}\,a_{t,ij}
    \right) \right\rangle .
\end{equation}

\begin{equation}\label{eq:algorithm_A01}
    \eta_{ij}=\frac{\mathcal{T}_{ij}}{\tau_{ij}},\qquad
    \mathcal A_{0,ij}=1-e^{-\eta_{ij}},\qquad
    \mathcal A_{1,ij}=1-(1+\eta_{ij})e^{-\eta_{ij}} .
\end{equation}

\begin{theorem}[Asymptotic-preserving continuum limit]\label{thm:ap}
Assume that the macroscopic variables remain smooth, that the interface GKS fluxes remain bounded, and that the time step is restricted by a mesh-scale CFL condition independent of the relaxation time. Let
\[
    \chi_i^\ast=\min\left(\chi_i,\min_{j\in N(i)}\eta_{ij}\right).
\]
In the limit $\chi_i^\ast\to\infty$, the WPD finite-volume update reduces to a Navier--Stokes GKS discretization with an exponentially small particle correction:
\begin{equation}\label{eq:ap_limit_update}
    \mathbf U_i^{n+1}
    =
    \mathbf U_i^n
    -
    \frac{\Delta t}{|\Omega_i|}
    \sum_{j\in N(i)}
    \left(
    \mathbf F_{\mathrm E,ij}^{\mathrm{GKS}}
    +
    \mathbf F_{\mathrm{vis},ij}^{\mathrm{GKS}}
    \right)|\Gamma_{ij}|
    +
    \mathcal O\!\left((1+\chi_i^\ast)e^{-\chi_i^\ast}\right) .
\end{equation}
\end{theorem}

\begin{proof}
From Eq.~\eqref{eq:algorithm_A01},
\begin{equation}
    \mathcal A_{0,ij}=1-e^{-\eta_{ij}}\to1,
    \qquad
    \mathcal A_{1,ij}=1-(1+\eta_{ij})e^{-\eta_{ij}}\to1
    \qquad
    \text{as }\eta_{ij}\to\infty .
\end{equation}
More precisely,
\[
    1-\mathcal A_{0,ij}=e^{-\eta_{ij}},
    \qquad
    1-\mathcal A_{1,ij}=(1+\eta_{ij})e^{-\eta_{ij}} .
\]
Therefore the wave flux in Eq.~\eqref{eq:wave_flux_algorithm} approaches the full Navier--Stokes GKS flux,
\begin{equation}
    \mathbf F_{\mathcal{W},ij}
    =
    \mathbf F_{\mathrm E,ij}^{\mathrm{GKS}}
    +
    \mathbf F_{\mathrm{vis},ij}^{\mathrm{GKS}}
    +
    \mathcal O\!\left((1+\eta_{ij})e^{-\eta_{ij}}\right) .
\end{equation}
The surviving particle memory is controlled by the factor $\beta_i=e^{-\chi_i}$ in the local-horizon decomposition and by the relaxation factor $e^{-\Delta t/\tau_i}$ during the particle update. For bounded particle moments and smooth reconstructed macroscopic states, its transported contribution is exponentially small in the continuum limit. Substituting these estimates into the conservative update \eqref{eq:fv_update_total} gives Eq.~\eqref{eq:ap_limit_update}. Since the time step is controlled by the acoustic CFL condition rather than by $\tau_i$, the scheme is asymptotic-preserving in the sense that it becomes a consistent Navier--Stokes finite-volume discretization without resolving the collision time.
\end{proof}

\end{document}
